% Generated draft: requires mathematical and source review.
\documentclass[11pt]{article}
\usepackage[margin=25mm]{geometry}
\usepackage{fontspec}
\setmainfont{Latin Modern Roman}
\usepackage{amsmath,amssymb,mathtools,unicode-math}
\setmathfont{Latin Modern Math}
\usepackage{enumitem,needspace,xurl,hyperref}
\setlength{\parindent}{0pt}
\setlength{\parskip}{5pt}
\setlength{\emergencystretch}{3em}
\newcommand{\N}{\mathbb{N}}
\newcommand{\Z}{\mathbb{Z}}
\newcommand{\Q}{\mathbb{Q}}
\newcommand{\R}{\mathbb{R}}
\newcommand{\C}{\mathbb{C}}
\newcommand{\F}{\mathbb{F}}
\newcommand{\A}{\mathbb{A}}
\newcommand{\PP}{\mathbb{P}}
\newcommand{\E}{\mathbb{E}}
\newcommand{\eps}{\varepsilon}
\newcommand{\dd}{\,\mathrm{d}}
\begin{document}
\section*{UnsolvedMath ID: 1100401}
\textbf{Source problem code:} catalogue entry
\subsection{Definitions and mathematical statement}
Identify the circle with the real projective line $S^1\cong\R\mathrm P^1$. Let
\[
 G=\operatorname{Homeo}_+(S^1)
\]
be the group of orientation-preserving circle homeomorphisms with the compact-open topology, equivalently uniform convergence for a compatible metric on the circle. Let
\[
 P=\operatorname{PSL}_2(\R)=\operatorname{SL}_2(\R)/\{I,-I\}
\]
be its usual subgroup acting projectively by $x\mapsto(ax+b)/(cx+d)$, with the usual interpretation at infinity.

A subgroup $P<G$ is maximal proper closed if it is closed and proper and every closed subgroup $K$ satisfying $P\subseteq K\subseteq G$ is either $P$ or $G$. De la Harpe's question is whether the standard projective subgroup $P$ is maximal proper closed in $\operatorname{Homeo}_+(S^1)$. Equivalently, does
\[
 \overline{\langle P,f\rangle}=\operatorname{Homeo}_+(S^1)
 \quad\text{hold for every }f\in\operatorname{Homeo}_+(S^1)\setminus P?
\]
The closure is in the compact-open topology.

The supplied assessment reports an affirmative resolution of this orientation-preserving question. It additionally says ``indeed of $\operatorname{Homeo}(S^1)$.'' That extra claim is incompatible with the stated maximality definition: $\operatorname{Homeo}_+(S^1)$ itself is a proper closed intermediate subgroup between $P$ and the full circle-homeomorphism group. The extra qualification is therefore retained as a source inconsistency requiring review, not incorporated into the mathematical assertion.
\subsection{Short English statement}
Is the projective action of $\operatorname{PSL}_2(\R)$ a maximal closed subgroup among orientation-preserving circle homeomorphisms? The supplied assessment attributes a yes answer to Giblin--Markovic (2006), resolving the question as stated. Its additional extension to the full homeomorphism group conflicts with the orientation-preserving intermediate subgroup and must be checked separately. The registry's solved label and the cited resolution are not independent verification in this draft.
\subsection{Sources}
\begin{itemize}
\item UnsolvedMath permanent ID 1100401, source code catalogue entry, ``Questions in Geometric Group Theory --- Q 4.1.'' Registry curation and metadata: Ulam AI and UnsolvedMath Contributors, CC BY 4.0; updated 2026-08-06.
\item Bestvina, \emph{Questions in Geometric Group Theory} (source list dated 2004), Question 4.1, attributed to de la Harpe, PDF page 12: \url{https://www.math.utah.edu/~bestvina/eprints/questions-updated.pdf}. Supplied manual transcription accessed 2026-07-29. Original-list metadata presents an open item and says current status was not independently verified there; status/release and rights review are NEEDS\_REVIEW.
\item Giblin--Markovic, \emph{Geometry \& Topology} \textbf{10} (2006), 1319--1346, Corollary in Section 1, as cited by the supplied resolving assessment. An article title or locator is not supplied; none is inferred.
\end{itemize}
\end{document}
